%% Chapter 1 — Introduction
\label{ch:intro}

\section{Motivation}
Many projects need specialized engineering skills. These skills cover software,
civil, electrical, mechanical, and network work. Projects keep getting more
complex, and more of them are now done remotely. So the demand for engineering
experts keeps growing.

Online service marketplaces have emerged as the primary response to this growing
demand, becoming the main way to find and hire independent professionals. The
online gig economy has grown steadily over the years~\cite{kassi2019online}.

Even so, finding the \textit{right} engineering expert is still slow and mostly
manual. General freelance platforms are built for high-volume, short,
transactional tasks, not for specialized engineering work such as CAD and MEP
drafting, embedded-firmware development, structural analysis, or network design.
Within these platforms, the fundamental challenge is discovery: the process by
which a client identifies the provider best suited to their specific project.

Discovery remains the bottleneck. Selection is manual and unstructured. A client
searches by keyword, then weighs skills, availability, cost, location, and
reputation by hand. There is no transparent, automatic way to rank providers
over these competing factors at once, and no way to say that, for a given job,
one factor matters more than another. This work gets harder as the number of
providers grows, and it often leads to a poor match.

Compounding this discovery challenge is the problem of trust. Existing platforms
do little technical vetting, allowing providers to overstate skills, and rely on
star ratings that are easy to inflate. This work does not attempt to fully
verify provider credentials, a task beyond its scope, but instead reduces the
opportunity for misrepresentation by design. Skills are drawn from a shared
catalogue with explicit proficiency levels rather than free text, ratings are
tied to completed engagements and awarded by both parties, and provider profiles
must pass a publication review before becoming discoverable.

Picking a provider by several factors at once is a well-known kind of problem.
The field of \textit{multi-criteria decision making} offers structured ways to
rank options that trade off competing factors~\cite{triantaphyllou2000multi}.
Using such a model for provider selection, instead of leaving it to manual
judgment, is the main idea of this work. The goal is an \textbf{intelligent
multi-criteria matching mechanism} that ranks engineering providers
automatically and clearly.

\section{Problem Statement}
The core problem is that current marketplaces have no clear, automatic way to
match clients with engineering providers. Existing solutions have three main
limits:

\begin{itemize}
	\item \textbf{Unstructured, manual discovery.} Provider selection relies on
	      keyword search and by-hand comparison. There is no transparent,
	      automatic way to rank providers over several competing factors at once.
	\item \textbf{No configurability.} Clients and operators cannot say that, for
	      a given job, one factor (such as cost) matters more than another (such
	      as location).
	\item \textbf{Limited engineering rigor.} Many platforms are not built or
	      evaluated for the secure communication, access control, and measured
	      network performance expected of a maintainable web service.
\end{itemize}

\section{Scope}
This work covers the design, build, and evaluation of a marketplace for
engineering services. It is organized around three core areas: foundation,
implementation, and validation.

\textbf{System Design and Architecture.} The work specifies functional
requirements and designs a layered client-server architecture (presentation,
business, and data tiers) with formal UML models. It defines the API contract
and documents all endpoints.

\textbf{Core Implementation.} The implementation encompasses a relational
database with strong data integrity, a maintainable backend service, and the
essential feature modules: user registration and authentication, role-based
access control for \textit{Client}, \textit{Service Provider}, and
\textit{Administrator} roles. Authentication and authorization are implemented
directly in the system (rather than delegated to external frameworks), ensuring
transparent control and enabling the security evaluation detailed in the
objectives. Additional features include service request management, provider
profile and skill management, and a bidirectional rating system. At the heart of
the system is an intelligent multi-criteria matching algorithm that ranks
providers across five weighted factors: skill compatibility, availability, cost
range, location proximity, and user rating, with adjustable weight assignment.
The work includes a comparison of basic keyword-matching and intelligent
multi-criteria approaches.

\textbf{Security and Evaluation.} The system enforces secure HTTPS
communication, implements session management, and applies defenses against
common web vulnerabilities. System evaluation covers functional correctness,
performance under varying user loads, the effectiveness of the matching
algorithm, and a discussion of scalability boundaries.

This pre-thesis report addresses the foundation and core implementation areas,
covering problem definition, related work, methodology, and system design. The
full matching engine implementation and the security and performance evaluation
are deferred to the later thesis phases outlined in the work plan.

\section{Live Deployment}
To aid review and testing, a live instance of the prototype is deployed and open
for exploration at \url{https://ems-portal-frontend.onrender.com/}. The
deployment uses a free tier that sleeps when idle, so the first request after a
quiet period can take up to a minute to wake; later requests are fast. Because
the entire stack is containerized with Docker, the same images that run locally
also run in production, making deployment straightforward.

Table~\ref{tab:demo} lists seeded demo accounts for testing the system with
different user roles. The demo database holds sample data only and may be reset
at any time without notice.

\begin{table}[ht]
	\centering
	\caption{Seeded demo accounts on the live deployment.}
	\label{tab:demo}
	\begin{tabular}{lll}
		\hline
		\textbf{Role} & \textbf{Email} & \textbf{Password} \\ \hline
		Client & \texttt{client@ems.local} & \texttt{password123} \\
		Service Provider & \texttt{provider@ems.local} & \texttt{password123} \\
		\hline
	\end{tabular}
\end{table}

\section{Objectives}
The objectives of this thesis are:

\begin{enumerate}
	\item To analyze the system requirements and design a scalable client-server
	      architecture for a service marketplace.
	\item To build a modular and maintainable web system that follows good
	      software engineering practice.
	\item To design and build an intelligent multi-criteria matching algorithm
	      that ranks and recommends providers.
	\item To build secure authentication and role-based authorization.
	\item To design and evaluate network communication efficiency and system
	      performance.
	\item To run functional testing, performance testing, and basic security
	      analysis.
\end{enumerate}

This pre-thesis report addresses objectives 1 to 4: requirement analysis and
architecture (1), modular system design and implementation (2), the design and
initial foundation of the matching algorithm (3), and secure authentication and
authorization (4).

The full matching engine and performance/security evaluation
(objectives 5 to 6) are carried out in the later thesis phases outlined in the work
plan.
